\section{A Dense but Weak Signal: Linear versus Nonlinear Models}
\label{sec:linvsnonlin}

Section~\ref{sec:marginal} established, one source at a time, that the exogenous
panel carries information beyond the HAR lag structure of \citet{Corsi2009}, and
that no single source carries much of it. That result is a statement about
\emph{features} under a fixed linear estimator. It leaves open the question this
section answers: how much of the achievable forecast improvement is a property of
the \emph{information set}, and how much is a property of the \emph{model class}
that reads it?

The organizing claim is that the exogenous signal here is \textbf{dense but
weak}. Many features each carry a small amount of predictive content, that
content is largely overlapping rather than complementary, and no small subset
concentrates it. This is the regime \citet{gu2020empirical} and
\citet{christensen2023machine} describe for asset-pricing and volatility panels
respectively, and it has four consequences documented in turn below. First, both
model classes extract real, statistically decisive skill --- neither fails.
Second, tree ensembles hold a modest but \emph{systematic} edge: they win on
every feature bucket we score, by a margin roughly a quarter the size of the
information effect itself. Third, that edge is not a hyperparameter-tuning story.
Fourth, the extra nonlinear capacity \emph{saturates almost immediately}:
additive terms plus shallow pairwise interactions recover essentially everything
a deep boosted ensemble finds.

\paragraph{A note on comparability of levels.}
QLIKE levels are not comparable across evaluation panels, and this section
necessarily touches four. Every table below is labeled with its panel and no
comparison is ever made across labels. They are: (i) the \textbf{causal-tuning
battery panel} --- the $n = 218{,}934$ out-of-sample panel of
Section~\ref{sec:descriptive}, per-bar refit, levels $\approx 0.126$--$0.135$;
(ii) the \textbf{rolling-tune panel} --- the same feature construction stitched
from 100 walk-forward chunks over the full 2005--2025 emit space, levels
$\approx 0.130$--$0.135$; (iii) the \textbf{deployed residual-stack panel} ---
the $n = 194{,}934$ panel on which the multi-stage residual models were
developed, levels $\approx 0.120$--$0.128$; and (iv) the older
\textbf{adjusted-scale model comparison} and the \textbf{deep-learning global
panel}, each with its own baseline and levels $\approx 0.11$--$0.48$. Within
each panel the ranking and the differences are the deliverable; the absolute
level is a property of the evaluation window and the warm-up convention.

\subsection{The battery: five estimators, nine buckets, one panel, one incumbent}
\label{sec:linvsnonlin_battery}

\paragraph{Design.}
The central experiment of this section is a fully crossed battery of five
estimators against nine exogenous feature buckets --- 45 arms, all complete.
Every arm is scored on the \emph{identical} out-of-sample panel of
$n = 218{,}934$ thirty-minute bars, with identical rows in every cell, and every
arm consumes the identical fixed feature construction; nothing varies across the
grid except the estimator and the bucket. The panel is computed as 100
walk-forward chunks over a rolling 500-trading-day ($24{,}000$-bar) training
window with a $24{,}000$-bar warm-up halo, pooled by concatenation, and the model
is refit \emph{every bar} throughout.

Critically, every arm re-selects its own hyperparameters periodically and
causally. Every 250 solves, the current training window is split forward into a
fit block, a 25-bar embargo (covering the 25-bar HAR moving average), and a
125-bar validation tail; the configuration minimizing validation MSE on that tail
is adopted. No forecast's hyperparameters ever see its own future, so the
comparison contains no look-ahead in the tuning channel --- the most common
source of inflated model-class effects in this literature. The machinery
implementing the sliding warm updates (exact rank-one linear algebra for the
ridge path, an exact every-bar online homotopy for the $\ell_1$ path) is
described in Section~\ref{sec:algorithms}.

The five estimators span the natural capacity ladder. Three are linear with
different shrinkage geometry: \textbf{rolling ridge} \citep{HoerlKennard1970},
\textbf{recursive elastic net} \citep{ZouHastie2005}, and \textbf{recursive
lasso} \citep{Tibshirani1996}, the last solved as an exact every-bar online
homotopy rather than a periodically re-run batch solver. Two are nonlinear
gradient-boosted tree ensembles: \textbf{LightGBM} \citep{Ke2017} and
\textbf{XGBoost} \citep{Chen2016}. Linear arms search log-scale penalty grids;
tree arms search max-depth $\{3,5,8\} \times$ colsample $\{0.25, 0.5\}$ with the
learning rate pinned at the production $0.1$, grid cells scored at 300 trees and
the winner refit every bar at the production 500 trees.

All 45 arms are compared against one shared comparator: the canonical incumbent
of Section~\ref{sec:descriptive}, per-bar OLS on HAR plus calendar features with
no exogenous inputs, at QLIKE $0.13415$ on this panel. Diebold--Mariano
statistics \citep{DieboldMariano1995} are computed on per-bar QLIKE loss
\citep{Patton2011} against that incumbent, with Newey--West HAC standard errors
and the small-sample correction of \citet{HarveyLeybourneNewbold1997}; a negative
$t$ favors the arm.

% Source: writeup/metrics_table_causal_tune_plus_spectral.csv (family == 'causal-tuned', 45 rows);
% cross-checked against writeup/causal_tune_battery_table.tex. QLIKE and |t| truncated to the
% precision of that table so the two agree cell-for-cell.
\begin{table}[H]
\centering
\small
\setlength{\tabcolsep}{4pt}
\caption{Causal-tuning battery: out-of-sample QLIKE by exogenous bucket and
estimator, with the Diebold--Mariano $t$-statistic against the shared OLS--HAR
incumbent in parentheses. \textbf{Panel: causal-tuning battery,}
$n = 218{,}934$ thirty-minute bars, identical rows in every cell; incumbent QLIKE
$= 0.13415$. Negative $t$ favors the arm; $|t| > 1.96$ is significant at the
5\% level. Best QLIKE in each row is bold. \texttt{baseline} is the no-exogenous
bucket: HAR plus calendar features only, so that row isolates the model class
with the information set held at the incumbent's.}
\label{tab:battery_bucket_estimator}
\begin{tabular}{@{}lccccc@{}}
\toprule
& \multicolumn{3}{c}{Penalized linear} & \multicolumn{2}{c}{Tree ensemble} \\
\cmidrule(lr){2-4} \cmidrule(lr){5-6}
Bucket & Ridge & Elastic net & Lasso & LightGBM & XGBoost \\
\midrule
\texttt{all\_features} & 0.12788 ($-19.9$) & 0.12827 ($-19.3$) & 0.12907 ($-13.2$) & \textbf{0.12604} ($-15.4$) & 0.12641 ($-16.5$) \\
\texttt{moments}       & 0.13085 ($-19.2$) & 0.13114 ($-17.7$) & 0.13179 ($-8.9$)  & \textbf{0.12816} ($-11.3$) & 0.12863 ($-9.0$) \\
\texttt{liquidity}     & 0.13120 ($-13.4$) & 0.13184 ($-11.5$) & 0.13206 ($-11.2$) & \textbf{0.12840} ($-14.7$) & 0.12886 ($-14.0$) \\
\texttt{market\_vw}    & 0.13291 ($-9.1$)  & 0.13316 ($-7.3$)  & 0.13321 ($-6.4$)  & \textbf{0.12963} ($-10.8$) & 0.13012 ($-8.1$) \\
\texttt{implied\_vol}  & 0.13252 ($-11.3$) & 0.13262 ($-10.6$) & 0.13274 ($-9.8$)  & \textbf{0.13001} ($-7.8$)  & 0.13077 ($-5.7$) \\
\texttt{market\_ew}    & 0.13319 ($-7.2$)  & 0.13344 ($-5.4$)  & 0.13353 ($-4.6$)  & \textbf{0.13039} ($-7.0$)  & 0.13043 ($-7.3$) \\
\texttt{sentiment}     & 0.13347 ($-5.7$)  & 0.13377 ($-3.5$)  & 0.13390 ($-2.1$)  & \textbf{0.13038} ($-8.1$)  & 0.13050 ($-8.4$) \\
\texttt{vol\_demand}   & 0.13337 ($-4.1$)  & 0.13361 ($-3.4$)  & 0.13359 ($-3.3$)  & \textbf{0.12979} ($-10.0$) & 0.13026 ($-8.6$) \\
\midrule
\texttt{baseline} (no exog) & 0.13445 ($+5.0$) & 0.13471 ($+7.5$) & 0.13475 ($+7.0$) & \textbf{0.13118} ($-5.6$) & 0.13163 ($-5.5$) \\
Incumbent: OLS on \texttt{baseline} & \multicolumn{5}{c}{0.13415} \\
\bottomrule
\end{tabular}
\end{table}

\paragraph{(a) The ordering is uniform, and the model-class effect is modest.}
Table~\ref{tab:battery_bucket_estimator} orders identically in every one of the
nine rows: tree ensembles beat penalized linear models, which beat the OLS--HAR
incumbent. LightGBM is the best arm in all nine buckets; ridge is the best linear
arm in all nine. There is no bucket on which the ranking reverses, and no bucket
on which the trees fail to find something the linear models cannot. This
uniformity is itself evidence for the dense-weak reading: if the nonlinear
advantage came from a few sharply interacting features it would be concentrated
in the buckets that contain them, and it is not.

The magnitude, however, is small. On the richest bucket, \texttt{all\_features},
the best nonlinear arm improves on the best linear arm by $0.12788 - 0.12604 =
0.00184$ QLIKE, while the total improvement of that best arm over the incumbent
is $0.13415 - 0.12604 = 0.00811$. The model class therefore accounts for
approximately $23\%$ of the recoverable gain, and the information set for the
remaining $77\%$. This ratio is the quantitative content of ``the lever is
information, not model,'' and it recurs at the same order of magnitude in every
subsequent experiment in this section.

\paragraph{(b) The linear gain is smaller but more consistent, exactly as a dense
weak signal predicts.}
The Diebold--Mariano statistics carry information the QLIKE levels do not, and on
\texttt{all\_features} they invert the ranking. Ridge, the \emph{worse} model by
QLIKE, achieves $t = -19.9$; LightGBM, the \emph{better} model, achieves only
$t = -15.4$; the elastic net shows the same pattern ($t = -19.3$ at QLIKE
$0.12827$). Because DM normalizes the mean loss differential by its HAC standard
error, a larger $|t|$ at a smaller mean improvement implies a \emph{smaller
bar-to-bar dispersion} of that improvement: the linear model wins by a little on
almost every bar, whereas the tree wins by more on average but with substantially
more variability, losing on a larger fraction of bars.

This is the signature of a dense weak signal read by two different mechanisms. A
penalized linear model aggregates a small, near-constant contribution from many
features, so its edge over HAR is a nearly deterministic offset with a
low-variance loss differential. A boosted tree instead localizes, finding regions
of the state space where a nonlinear correction is worth making and leaving the
rest to its additive base; localized corrections are larger where they apply and
absent where they do not, raising the mean improvement and its variance together.
Neither behavior is pathological and the two are complementary, but the inversion
is a useful warning that a QLIKE ranking and a significance ranking need not
agree.

\paragraph{(c) The only losses to plain OLS isolate the pure-nonlinearity
premium.}
Of the 45 arms, exactly three are worse than the shared incumbent, and they are
the three penalized linear arms on the no-exogenous \texttt{baseline} bucket:
ridge $0.13445$ ($t = +5.0$), elastic net $0.13471$ ($t = +7.5$), lasso
$0.13475$ ($t = +7.0$), all decisively significant. This is a useful negative
control on the whole shrinkage apparatus. When the only regressors present are
the six HAR lags and a handful of calendar dummies --- a low-dimensional,
well-conditioned, entirely informative design --- shrinking them strictly hurts.
Shrinkage is a variance-reduction device, and here there is no estimation
variance worth reducing, so the bias it induces is pure cost. Penalized linear
models earn their keep only once the design matrix contains a large number of
weak or dead columns, which is precisely what the exogenous buckets supply.

The trees on that same \texttt{baseline} bucket do the opposite. LightGBM
improves the incumbent from $0.13415$ to $0.13118$, a gain of $0.00297$ with
$t = -5.6$; XGBoost reaches $0.13163$ ($t = -5.5$). Because the information set
in that row is \emph{identical} to the incumbent's --- same HAR lags, same
calendar features, no exogenous inputs whatsoever --- this $0.00297$ is a clean
measurement of the \textbf{pure-nonlinearity premium}: the value of relaxing the
linear functional form on HAR and calendar features alone. Set beside the
\texttt{all\_features} row it yields the section's cleanest decomposition: of the
$0.00811$ total gain of the best arm over the incumbent, roughly $0.00297$ is
available from nonlinearity in the HAR-plus-calendar block by itself, and the
remainder comes from the exogenous information. The two premia are of the same
order and largely additive, with the information side the larger.

\paragraph{(d) Limitation.}
The DM statistics in Table~\ref{tab:battery_bucket_estimator} are computed
against the shared OLS--HAR incumbent \emph{only}. They establish that each arm
differs significantly from the incumbent; they do \emph{not} establish that any
arm differs significantly from any other arm. Pairwise model-versus-model DM
tests and a Model Confidence Set \citep{HansenLundeNason2011} over all 45 arms
are computed but not yet complete at the time of writing. Consequently the
$0.00184$ linear-to-nonlinear gap, and the uniform LightGBM-over-ridge ordering
that supports it, carry \emph{no significance statement}. They are a consistent
point-estimate pattern across nine independent buckets, which is suggestive, and
nothing stronger. We flag this explicitly because the effect sizes at issue
($10^{-3}$ QLIKE) are of the same order as several effects elsewhere in this
paper that \emph{did} fail formal significance testing.

\subsection{The edge is not a tuning story}
\label{sec:linvsnonlin_tuning}

A natural alternative explanation for Table~\ref{tab:battery_bucket_estimator} is
effort: tree ensembles have more hyperparameters, so a tuned tree is simply a
harder-tried model than a tuned linear one, and the gap measures search budget
rather than function class. Three experiments address this, and they point in a
more interesting direction than either a clean confirmation or a clean refutation.

\paragraph{(a) Causal hyperparameter tuning of the linear penalty is worse than
doing nothing.}
The sharpest test is to ask whether \emph{any} leakage-free re-selection of the
linear penalty improves on an untouched library default. The experiment holds
features, executor path, and refit cadence fixed and varies only the penalty
channel. The tuned arm runs a single causal selector over a pooled grid --- six
ridge penalties on an exact rank-one path, five lasso penalties, and five
elastic-net penalties on a warm online homotopy --- re-selecting every 250 bars
from a forward fit/embargo/validation split of the trailing window, with
coefficients refit every bar. The comparator is scikit-learn's own Ridge default
$\alpha = 1.0$, per-bar refit, with no selection at all.

% Source: writeup/rolling_tune_meeting_2026-07-09.md, "Linear results" table.
\begin{table}[H]
\centering
\small
\caption{Causal rolling hyperparameter selection versus an untouched library
default, by exogenous bucket. \textbf{Panel: rolling-tune,} full 2005--2025 emit
space stitched from 100 walk-forward chunks. Positive $\Delta$ means the causally
tuned arm is \emph{worse}. The \texttt{sentiment} and \texttt{vol\_demand} rows
land 97 and 95 of 100 chunks respectively and are reported net of a localized
solver pathology at the exogenous data-start boundary (see text); all other rows
are complete.}
\label{tab:causal_tune_vs_default}
\begin{tabular}{@{}lrrrr@{}}
\toprule
Bucket & Chunks & Causally tuned & Library default & $\Delta$ (tuned $-$ default) \\
\midrule
\texttt{moments}      & 100/100 & 0.131459 & 0.130264 & $+0.001194$ \\
\texttt{liquidity}    & 100/100 & 0.131566 & 0.130822 & $+0.000744$ \\
\texttt{implied\_vol} & 100/100 & 0.132856 & 0.132091 & $+0.000765$ \\
\texttt{market\_vw}   & 100/100 & 0.133015 & 0.132388 & $+0.000627$ \\
\texttt{market\_ew}   & 100/100 & 0.133233 & 0.132677 & $+0.000556$ \\
\texttt{sentiment}    & 97/100  & 0.134395 & 0.133581 & $+0.000814$ \\
\texttt{vol\_demand}  & 95/100  & 0.134480 & 0.133754 & $+0.000726$ \\
\bottomrule
\end{tabular}
\end{table}

The result is uniform and one-signed: the leakage-free tuner \emph{trails} the
untouched default by $+0.0006$ to $+0.0012$ QLIKE on all seven landed buckets.
Causal selection does not merely fail to help; it pays a small, consistent
variance tax. Every selection is made from a 125-bar validation tail, which at
this signal-to-noise ratio is a noisy statistic, and acting on that noise every
250 bars injects more variance than the penalty choice can recover. The
selector's own behavior confirms the landscape is flat: across all tunings it
never collapses onto a single penalty family, choosing ridge 41--50\% of the
time, lasso 29--35\%, and elastic net 20--25\%, in every bucket. A selector
facing a genuinely differentiated objective concentrates; one facing three
near-equivalent options distributes according to validation noise. (For
completeness: 8 of 700 chunks diverged, all at the exogenous data-start boundary
around 2010--2012 where columns constant for years switch on mid-window. Under
forward-fill plus availability indicators the Gram matrix passes through
near-singularity, and the pure-lasso grid arm --- $\ell_2$ weight exactly zero,
the one solver in the grid with no well-posedness guarantee --- is the only arm
affected; the default arm was never affected.)

\paragraph{(b) Even an oracle allowed to cheat cannot buy more.}
The rolling-tune result establishes that \emph{causal} selection is worthless. A
stronger question is whether \emph{any} penalty configuration, chosen with full
knowledge of the test set, would have been materially better. We answer it by
running Optuna's TPE sampler \citep{Akiba2019} with the full out-of-sample QLIKE
as the objective directly. This is test-set overfitting by construction and the
resulting numbers are \emph{not} out-of-sample performance; they are a diagnostic
upper bound on what the penalty channel could ever deliver.

Against the incumbent elastic net on the same base ($0.12540$ on that panel), the
oracle finds: lasso best $0.12524$, an improvement of $-0.00015$, with the
penalty at the search-grid floor; elastic net best $0.12529$ ($-0.00011$); and
ridge best $0.12578$, which is $+0.00038$ \emph{worse} than the incumbent ---
ridge loses even as an oracle. The ceiling on the entire penalty channel, with
cheating permitted, is therefore $-0.00015$ QLIKE, roughly an order of magnitude
below the $0.00184$ model-class gap and comfortably inside walk-forward noise. A
proper cross-validated procedure would capture some fraction of that ceiling,
which is to say nothing at all. Every oracle winner also sits at minimal
shrinkage, so the ceiling is attained by approaching the unpenalized solution.

That ridge loses even as an oracle on the dense feature set has a concrete
mechanism, and it is why the elastic net rather than ridge anchors the production
linear path. The lasso support analysis shows $304$ of $523$ exogenous columns
zeroed, of which $235$ --- $77\%$ --- are dead or constant availability
indicators. $\ell_2$ shrinkage can make those columns small but cannot make them
zero, so they continue to inject estimation variance; $\ell_1$ removes them
outright. Where a bucket contains no such junk the relationship inverts: on
\texttt{implied\_vol}, a dense block of all-informative VIX levels, ridge is the
best penalty, because zeroing correlated informative columns discards signal.
Penalty choice tracks the junk fraction of the design and nothing else --- which
is another way of saying it is not a lever.

\paragraph{(c) The flip side: defaults \emph{do} mislead, but only for trees.}
The preceding two results should not be read as a general claim that
hyperparameters do not matter. They matter enormously in one place: the gap
between a tree ensemble's library defaults and a properly configured one. An
earlier model comparison on this project, run on the \textbf{adjusted-scale
model-comparison panel} --- a different panel with a different baseline and
levels around $0.11$--$0.35$, \emph{not} comparable to any other table in this
section --- makes the point starkly. On HAR-only features, ridge scores QLIKE
$0.1380$. Default-hyperparameter LightGBM scores $0.2401$ and default-hyperparameter
XGBoost $0.2407$: the trees lose to ridge by approximately $0.10$ QLIKE units, an
enormous margin that would comfortably support the conclusion that linear HAR is
uniquely suited to realized-volatility forecasting. After Bayesian tuning, that
conclusion evaporates. Tuned LightGBM reaches $0.1099$ --- beating ridge outright
by $0.028$ --- and tuned XGBoost $0.1264$, placing every tuned tree within
$0.005$ QLIKE of ridge or ahead of it.

The failure of the defaults is legible in the hyperparameter deltas: one hundred
depth-three trees are severely capacity-constrained on $\approx 218{,}000$
observations with a dense multi-scale lag basis, and the search fixes capacity
and regularization simultaneously (depth $3 \to 6$ and $100 \to 300$ estimators
for XGBoost; $31 \to 59$ leaves and $100 \to 1000$ estimators for LightGBM;
learning rate roughly halved for both; row and column subsampling and explicit
penalties switched on). The methodological implication generalizes beyond this
dataset and connects to the framing in \citet{gu2020empirical} and
\citet{christensen2023machine}: the widely repeated observation that tree
ensembles lag linear models on realized-volatility tasks is, at least in part, a
\textbf{default-hyperparameter artifact}. The effort gap between a closed-form
linear estimator, which has essentially nothing to configure, and an
out-of-the-box boosted ensemble is large enough to masquerade as a model-class
effect; any comparison of the two without a stated tuning protocol should be
treated as uninformative.

The two halves of this subsection are consistent rather than contradictory.
Configuring a tree \emph{into} the right capacity regime is worth
$\approx 0.10$ QLIKE on that panel; once it is there, further causal re-tuning is
worth nothing. (The tree half of the rolling-tune experiment returned
$0.176807$ tuned versus $0.177566$ default on its single landed canary chunk ---
opposite in sign to the linear result, but one chunk is noise and we report it as
such rather than as a finding.) Hyperparameters are a threshold to be cleared,
not a dimension to be optimized along.

\subsection{Where the nonlinearity lives, and where it saturates}
\label{sec:linvsnonlin_saturation}

Granting that the trees hold a real edge, what is it made of? Three independent
decompositions agree that it is shallow.

\paragraph{(a) Functional ANOVA: mostly additive, with a diffuse remainder.}
Decomposing the fitted depth-eight boosted tree into orders of interaction gives
approximately $55\%$ additive, $9\%$ pairwise, and $36\%$
higher-order-and-diffuse. Within the additive part, HAR dominates: the HAR block
alone accounts for $96\%$ of the explainable variance of the full fit. The $36\%$
higher-order share is not concentrated in a few strong interactions --- it is
spread thinly across many, and when individual pairs are tested against a placebo
null only the VIX-term interactions survive. This is a model nominally capable of
eighth-order interactions that in practice spends most of its capacity on main
effects.

The same conclusion arrives from the feature side. The signal is \emph{dense},
not low-rank: the elastic-net solution retains approximately $120$ effective
dimensions out of $529$ features, and principal-component truncation fails
outright --- five components retain $72\%$ of the feature variance while losing
essentially all of the predictive signal. $\ell_1$ selection in the raw basis
works where variance-directed rotation does not, so the signal is \emph{sparse in
the raw basis but not low-rank in covariance}. That distinction is precisely why
ridge-style dense shrinkage underperforms on the full panel, and why no
dimension-reduction preprocessing step rescues it.

\paragraph{(b) The expressivity ladder: pairwise interactions are the whole
story.}
The most direct test of saturation is to hold the base model fixed and vary only
the expressivity of a correction stage fitted to its residual. The following
ladder does exactly that on the \textbf{deployed residual-stack panel}
($n = 194{,}934$; levels $\approx 0.120$--$0.128$; not comparable to
Tables~\ref{tab:battery_bucket_estimator} or~\ref{tab:causal_tune_vs_default}).
The reference points are the base with no correction stage at all, $0.12081$, and
the deployed bagged additive-plus-pairwise stage, $0.12033$.

% Source: writeup/regime_moe_results_2026-06-29.md, "The expressivity ladder" table
% (sat_d0 / fm_d0_wd10 / bfm_d0_pure rows) and the ctrl_ebm anchor; XGB regime figure
% from writeup/mechanism_and_data_to_buy_2026-06-28.md (120-trial async-Optuna tune).
\begin{table}[H]
\centering
\small
\caption{Expressivity ladder for the residual correction stage.
\textbf{Panel: deployed residual-stack,} $n = 194{,}934$, full out-of-sample.
Each row replaces only the correction stage; the base model is identical
throughout. ``No correction stage'' is the base alone.}
\label{tab:expressivity_ladder}
\begin{tabular}{@{}llr@{}}
\toprule
Correction stage & Function class & QLIKE \\
\midrule
Additive neural GAM              & main effects only                    & 0.12757 \\
Bilinear factorization machine   & $+$ linear pairwise interactions      & 0.12108 \\
Hinge-basis FM, pure pairwise    & $+$ nonlinear shapes and interactions & 0.12080 \\
\midrule
\emph{No correction stage}       & ---                                   & \emph{0.12081} \\
Optuna-tuned XGBoost stage       & unrestricted depth                    & 0.12094 \\
\midrule
Bagged additive $+$ pairwise (GA2M) & main effects $+$ pairwise, bagged  & \textbf{0.12033} \\
\bottomrule
\end{tabular}
\end{table}

Two readings follow. First, \textbf{each expressivity upgrade closes the gap
monotonically, and the closing is essentially complete at pairwise}. A purely
additive stage is catastrophically inadequate at $0.12757$ --- worse than fitting
no correction at all. Adding bilinear interactions recovers almost the entire
deficit in one step, to $0.12108$; allowing those interactions nonlinear shapes
adds a further $0.00028$, reaching $0.12080$ and exactly matching the
no-correction baseline. The deployed stage, a bagged GA2M in the sense of
\citet{Lou2013} --- additive shape functions plus explicitly enumerated pairwise
terms --- is best of all at $0.12033$. Its remaining advantage over the
hinge-basis FM is attributable to bagging, not expressivity: the stage is fitted
on a few thousand rows of an approximately $93\%$-noise residual, where variance
control rather than capacity binds. That an un-bagged flexible fit \emph{hurts}
on this sample while a bagged additive-plus-pairwise fit helps is the clearest
available statement that we are at a variance floor, not a bias floor.

Second, and decisively, \textbf{unrestricted capacity is worse than restricted
capacity}. A 120-trial asynchronous Optuna search over the XGBoost correction
stage, warm-started from the best hand-curated configuration, lands at $0.12094$
--- worse than the interpretable GA2M at $0.12033$, and worse than fitting
\emph{no} correction stage at all at $0.12081$. Properly searched, unrestricted
depth on this residual is a net negative. The GA2M is therefore not an
interpretability compromise purchased at a cost in accuracy; here it is
simultaneously the more interpretable and the more accurate choice.

\paragraph{(c) The weak-signal floor, and one line that captures most of it.}
Two measurements bound how much structure remains. First, an \emph{in-sample}
GA2M probe of the leftover residual --- in-sample, hence an upper bound on what
any model could extract --- achieves $R^2 \approx 0.07$. Roughly $93\%$ of what
remains is irreducible noise, and any model-class comparison conducted in that
regime is a comparison of small differences in how gracefully each class handles
variance.

Second, and more constructively, the nonlinearity that does matter is
substantially recoverable in closed form. On an earlier build of the same
residual-stack panel family, a plain elastic-net base scores $0.12516$ and a
residualized boosted-tree correction takes it to $0.12414$. Adding a \emph{single
hand-written linear interaction} --- HAR lags multiplied by a late-day session
dummy, six features of which $\ell_1$ retains about four --- reaches $0.12457$,
recovering $58\%$ of the tree's entire edge. Expanding to a full set of
HAR-by-six-regime interactions (36 features) reaches $0.12453$, or $62\%$: the
marginal regimes beyond late-day add $4\%$ between them.

This distillation identifies what the tree is doing. It is not discovering a rich
high-order manifold; it is discovering that HAR's volatility-persistence
coefficient changes sign near the session close, and encoding that change through
splits. The mechanism is \emph{clock-anchored} rather than state-anchored, which
Section~\ref{sec:descriptive} documents directly in the raw data as a sharp
sign-flip in the correlation between short-horizon realized volatility and the
next-bar residual at the open and at the close/after-hours transition. Once the
regime is named and written as a product term, most of the nonlinear advantage
transfers to the linear model; the residual $\approx 38\%$ is what the tree still
earns, and Table~\ref{tab:expressivity_ladder} says it is pairwise.

\paragraph{(d) The stage ladder.}
Assembling the pieces in the order they were added gives the deployed stack, and
a compact summary of the whole model-side program.

% Source: writeup/mechanism_and_data_to_buy_2026-06-28.md, section 1 "Bottom line" table;
% final two rows from writeup/regime_moe_results_2026-06-29.md (EBM anchor 0.12033) and
% writeup/sections/meeting_update_2026_06.tex (composed floor 0.12022).
\begin{table}[H]
\centering
\small
\caption{Cumulative stage ladder of the deployed residual stack.
\textbf{Panel: deployed residual-stack,} full out-of-sample, raw-variance scale
after Duan smearing. Each row adds one stage to the row above.}
\label{tab:stage_ladder}
\begin{tabular}{@{}lrr@{}}
\toprule
Stage & QLIKE & $\Delta$ \\
\midrule
Plain elastic net                                & 0.12516 & --- \\
$+$ engineered regime and path features          & 0.12314 & $-0.00202$ \\
$+$ global depth-8 XGBoost on the residual       & 0.12129 & $-0.00185$ \\
$+$ GA2M correction stage on the session-edge window & 0.12050 & $-0.00079$ \\
Best composed stack                              & 0.12033 & $-0.00017$ \\
\midrule
Floor with ensemble-weight tuning                & 0.12022 & $-0.00011$ \\
\bottomrule
\end{tabular}
\end{table}

The shape of the ladder repeats the section's thesis at the level of the whole
pipeline. The two largest steps --- engineered features and the first nonlinear
stage --- are of comparable size, $-0.00202$ and $-0.00185$; everything after them
is fifth-decimal. The final $-0.00011$, obtained by tuning the weight of an
ensemble mixture, sits below any threshold we are prepared to treat as real on a
panel of this length, and we report it as the floor rather than as a gain.

\paragraph{(e) Deep sequence models fail outright at this scale.}
For completeness we report the one model class that does not merely saturate but
fails. PatchTST \citep{Nie2023}, a patch-based transformer for multivariate time
series, was trained as a global model on its own evaluation panel
(\textbf{panel: deep-learning global}; baseline QLIKE $0.25145$ at
$n = 212{,}364$, levels not comparable to any other table here). Its best
configuration reaches QLIKE $0.46125$ at $n = 221{,}579$ --- roughly twice the
loss of its own panel baseline --- with an out-of-sample $R^2$ of $-0.77$ against
that baseline; a second configuration using overlapping patches is worse still at
$0.47856$. This is not evidence that transformers cannot forecast realized
variance. It is evidence that at $\approx 200{,}000$ observations of a signal
this weak, a model class whose sample-efficiency argument depends on abundant
data has no path to competitiveness, and that the honest reading of a $-0.77$
out-of-sample $R^2$ is failure rather than a tuning deficit to be excused. The
result is consistent with everything else in this section: capacity is not the
binding constraint, and adding more of it is actively harmful.

\subsection{Synthesis}
\label{sec:linvsnonlin_synthesis}

The dense-weak regime has a practical law, and this section has measured all
three of its terms on a single panel:
\[
\text{information} \;\gg\; \text{model class} \;\gg\; \text{hyperparameters}.
\]
The \textbf{information set} is worth $0.006$--$0.008$ QLIKE: handing the
incumbent the exogenous panel moves it from $0.13415$ to $0.12788$ under a linear
model and to $0.12604$ under a tree. The \textbf{model class} is worth
$0.002$--$0.003$: the pure-nonlinearity premium measured on the no-exogenous
bucket is $0.00297$, and the best-linear-to-best-nonlinear gap on the full
feature set is $0.00184$. \textbf{Hyperparameters}, once a model is in the right
capacity regime, are worth approximately zero: causal re-tuning of the linear
penalty \emph{costs} $+0.0006$ to $+0.0012$ uniformly, and an oracle permitted to
overfit the test set caps out at $-0.00015$. The one exception --- clearing a
tree out of its library defaults, worth $\approx 0.10$ QLIKE on the older panel
--- is a threshold, not a gradient, and once cleared it never has to be revisited.

Both model classes work and both should be reported. Linear models with
$\ell_1$-capable shrinkage win on consistency --- the larger Diebold--Mariano
statistics at the smaller mean improvement --- and remain the right base layer
for a signal spread thinly across hundreds of columns of which several hundred
are dead. Tree ensembles win on level, uniformly and on every bucket, by
localizing a correction the linear model cannot express. But that correction is
shallow: additive terms plus pairwise interactions recover essentially all of it;
unrestricted depth, properly searched, does worse than no correction stage at
all; and a single hand-written product of HAR with a late-day dummy recovers
$58\%$ of the tree's entire advantage in closed form. What the nonlinearity is
\emph{about} is not a rich functional geometry but a clock-anchored regime at the
session boundaries --- which is why it distills so cleanly, and why, having
distilled it, further capacity finds nothing.

The forward implication for Section~\ref{sec:algorithms} is that the interesting
algorithmic questions here are not questions of capacity, but of how to fit a
weak, drifting, high-dimensional signal \emph{exactly and cheaply enough} to
refit at every bar over a $218{,}934$-bar panel --- rank-one linear algebra,
online homotopy paths, and leakage-clean causal selection.
